\documentclass[11pt]{article}
\usepackage[utf8]{inputenc}
\usepackage[T1]{fontenc}
\usepackage{amsmath,amssymb}
\usepackage{graphicx}
\usepackage{hyperref}
\usepackage{booktabs}
\usepackage{url}
\usepackage[margin=1in]{geometry}

\title{M31 Autonomous: A Terminal-Native AI Coding Agent with Eight-Phase Workflow Orchestration}

\author{
  Eshan Roy\\
  Independent Researcher\\
  \texttt{eshanized@proton.me}\\
  \href{https://orcid.org/0009-0007-1261-6805}{ORCID: 0009-0007-1261-6805}
}

\date{June 22, 2026}

\begin{document}

\maketitle

\begin{abstract}
M31 Autonomous is a terminal-native AI coding agent, written entirely in Go, that orchestrates an eight-phase software engineering workflow end-to-end. From project initialization through discussion, planning, execution, verification, runtime smoke testing, and shipping, every run concludes with a verified git commit and a cross-session learning record. The system is distributed as a static binary (15--20 MB) with zero runtime dependencies and zero telemetry.

The architecture follows a strict layered design: a 33-screen terminal UI built on Bubble Tea, an eight-phase workflow engine, a provider abstraction layer supporting multiple LLM backends with automatic fallback, a permission-gated tool system with 17 registered tools, a composable skills framework, and a concurrent session coordinator with demand coalescing.

Key innovations include: (1) a complexity-adaptive workflow engine that classifies user goals into trivial, simple, moderate, or complex categories and selects an appropriate phase subset; (2) an AutoDream context consolidation system; (3) an automatic session compaction engine; (4) a runtime verification phase with HTTP smoke testing; and (5) a cross-session learning ledger. The system's codebase intelligence subsystem (EFIE) implements novel algorithms for community-structured, importance-weighted codebase exploration with formal mathematical proofs of correctness and complexity bounds.
\end{abstract}

\section{Summary}

M31 Autonomous is a terminal-native AI coding agent, written entirely in Go, that orchestrates an eight-phase software engineering workflow end-to-end. Beyond its practical utility as a development tool, M31A serves as a research platform for computational software engineering, hosting novel algorithms with formal mathematical foundations and enabling reproducible research workflows through its built-in research phase and cross-session learning ledger.

The architecture follows a strict layered design: a 33-screen terminal UI built on Bubble Tea~\cite{bubbletea}, an eight-phase workflow engine, a provider abstraction layer supporting multiple LLM backends with automatic fallback, a permission-gated tool system with 17 registered tools, a composable skills framework, and a concurrent session coordinator with demand coalescing.

Key innovations include: (1) a complexity-adaptive workflow engine that classifies user goals into trivial, simple, moderate, or complex categories and selects an appropriate phase subset, avoiding unnecessary overhead for simple tasks; (2) an AutoDream context consolidation system that compresses older messages into compact memory segments while protecting system prompts, tool calls, and recent context; (3) an automatic session compaction engine that triggers LLM-driven summarization when token usage approaches the model's context window limit; (4) a runtime verification phase that starts a development server, discovers routes from goal text and filesystem conventions, and runs concurrent HTTP smoke tests to catch errors that static checks miss; and (5) a cross-session learning ledger that records goal, task outcomes, file patterns, model used, and duration for every shipped session, enabling the agent to learn from its own history.

The workflow engine implements phase transition guards and a plan-discuss oscillation guard (capped at three cycles) to prevent infinite refinement loops. The plan parser uses cascading regex strategies with fallback JSON extraction, validated by duplicate detection, cycle detection via iterative DFS, and granularity checks. The execute phase uses Kahn's algorithm~\cite{kahn1962} for topological task scheduling with bounded concurrency and per-task self-healing loops that fall back to git bisect when standard healing fails.

\section{Statement of Need}

AI-assisted coding tools fall into two paradigms: editor-embedded assistants that sacrifice terminal flexibility, and command-line wrappers around single LLM calls that lack structured workflow. Real engineering tasks such as migrating an authentication middleware require understanding the codebase, discussing tradeoffs, creating a plan, executing changes across files, verifying correctness, and committing the result. Each phase demands different capabilities, autonomy levels, and safety guarantees.

From a research perspective, M31A addresses the gap between practical software engineering tools and formal computational methods. The system's codebase intelligence subsystem (EFIE) demonstrates how graph-theoretic algorithms---PageRank, Louvain community detection, and betweenness centrality---can be applied to code exploration with provable complexity bounds. This makes M31A both a useful tool and a platform for studying AI-assisted software engineering.

M31 Autonomous addresses this gap by treating the terminal as the primary interface and owning the entire task lifecycle. Unlike existing tools such as Aider~\cite{aider} and Cline~\cite{cline}, which provide conversational code editing with git integration, M31 Autonomous provides a structured workflow engine with explicit phases, automatic complexity classification, runtime verification with smoke testing, cross-session learning, and composable user-defined skills.

\section{Software Design}

M31 Autonomous follows a strict layered architecture with a clear dependency rule: higher layers may depend on lower layers, but never the reverse. The system comprises six major subsystems:

\begin{enumerate}
    \item \textbf{Terminal User Interface:} A 33-screen TUI built on Bubble Tea using the Elm architecture, providing keyboard-driven efficiency with screen transitions, command palette, and theming.
    
    \item \textbf{Workflow Engine:} An eight-phase pipeline (Initialize, Discuss, Plan, Execute, Verify, Runtime, Ship) with complexity-adaptive mode selection.
    
    \item \textbf{Provider Abstraction:} A multi-provider LLM layer supporting OpenRouter and Zen with automatic fallback, health checks, and cost tracking via atomic CAS operations.
    
    \item \textbf{Tool System:} 17 registered tools with permission-gated execution, rate limiting, and risk-level classification.
    
    \item \textbf{Code Intelligence (EFIE):} A novel community-structured, importance-weighted codebase exploration algorithm with formal mathematical foundations, including PageRank-based importance scoring, Louvain community detection, and betweenness centrality approximation.
    
    \item \textbf{Cross-Session Learning:} A persistent ledger recording session outcomes, enabling the agent to learn from its own history across sessions.
\end{enumerate}

\section{State of the Field}

Terminal-native AI coding agents have emerged as a distinct category following the adoption of large language models for code generation~\cite{chen2021codex}. Aider~\cite{aider} provides git-integrated multi-file editing with support for multiple LLM providers but does not include a structured workflow engine, runtime verification, or cross-session learning. Cline~\cite{cline} operates as a VS Code extension with tool-use capabilities but is not terminal-native. OpenHands~\cite{openhands} and SWE-agent~\cite{sweagent} target autonomous issue resolution on benchmarks rather than interactive developer workflows.

M31 Autonomous distinguishes itself through its eight-phase workflow engine with complexity-adaptive mode selection, runtime smoke testing, composable skills, and persistent cross-session memory.

\section{Research Impact Statement}

M31 Autonomous serves as both a software engineering tool and a platform for computational research. The system's codebase intelligence subsystem (EFIE) implements novel algorithms for community-structured, importance-weighted codebase exploration, with formal mathematical proofs of correctness, convergence guarantees, and complexity bounds published in the repository's \texttt{EFIE/} directory.

Key research contributions include:

\begin{enumerate}
    \item \textbf{Novel Algorithm Design:} The EFIE algorithm introduces adaptive expansion with community detection (Louvain modularity)~\cite{blondel2008}, PageRank-based importance scoring~\cite{brin1998}, and betweenness centrality approximation---a combination not previously applied to codebase exploration.
    
    \item \textbf{Formal Mathematical Rigor:} The EFIE research paper contains 11 formal definitions, 3 numbered theorems with proofs, and complete complexity analysis ($O(N \times F / P + E \times \log V)$ build time, $O(S \times B)$ query time).
    
    \item \textbf{Reproducible Research Workflows:} M31A's built-in research phase supports systematic codebase investigation before implementation planning, enabling reproducible software engineering research.
    
    \item \textbf{Cross-Session Learning:} The cross-session learning ledger records goal, task outcomes, file patterns, model used, and duration for every session, creating a persistent record suitable for longitudinal research studies.
\end{enumerate}

\section{AI Usage Disclosure}

No generative AI tools were used in the development of this software. The EFIE algorithm, its mathematical proofs, and the M31 Autonomous workflow engine were designed and implemented by the author without AI assistance.

The paper manuscript was written by the author. No AI writing assistants were used in drafting, editing, or revising this paper.

\section*{Acknowledgements}

The author thanks the open-source community for their contributions to the Bubble Tea framework and the broader Go ecosystem.

\bibliographystyle{plain}
\bibliography{paper}

\end{document}
